\ifdefined\CRevisionRound\else\def\CRevisionRound{2}\fi
\documentclass[10pt]{article}
\pdfvariable suppressoptionalinfo 767
\usepackage[a4paper,margin=22mm]{geometry}
\usepackage{amsmath,amssymb,amsthm,booktabs}
\usepackage{fontspec}
\setmainfont{TeX Gyre Pagella}
\setmonofont{Latin Modern Mono}
\usepackage{unicode-math}
\setmathfont{TeX Gyre Pagella Math}
\usepackage[english]{babel}
\babelprovide[import]{chinese-simplified}
\babelfont[chinese-simplified]{rm}[Renderer=Harfbuzz,Language=Default]{Droid Sans Fallback}
\usepackage[hidelinks]{hyperref}
\newtheorem{theorem}{Theorem}
\newtheorem{lemma}{Lemma}
\newtheorem{corollary}{Corollary}
\newcommand{\ii}{\mathrm i}
\newcommand{\ee}{\mathrm e}
\newcommand{\ZZ}{\mathbb Z}
\newcommand{\RR}{\mathbb R}
\newcommand{\dd}{\,\mathrm d}
\newcommand{\tr}{\operatorname{tr}}
\ifcase\CRevisionRound
\title{Reversible Multibaker Transport:\\A Complete Boundary-Aware Primitive Orbit Atlas}
\or
\title{Reversible Multibaker Transport:\\Primitive Orbits and the Corrected Twisted Zeta}
\else
\title{Reversible Multibaker Transport:\\Boundary-Corrected Zeta, Diffusion, and Parity-Sensitive Relaxation}
\fi
\author{HCS-C379 theorem and reproducibility package}
\date{5 September 2026}
\begin{document}
\maketitle
\begin{abstract}
We give a complete source-system theorem for the reversible multibaker on a
finite ring, with a domain convention that resolves the binary endpoint
ambiguity. On the invariant set of non-dyadic coordinates, every geometric
periodic point is reconstructed from a mixed binary word. A primitive
necklace of length $d$ and displacement $S$ lifts to $\gcd(L,S)$ cycles of
least period $dL/\gcd(L,S)$, with exact winding, reversal and reciprocal
stability multipliers.
\ifnum\CRevisionRound>0
The finite tilted transport determinant has a Chebyshev formula. Its weighted
primitive product contains two homogeneous symbolic cycles absent from the
geometric domain; removing them gives an explicit all-order correction.
We distinguish inverse-unstable weights from two-dimensional flat-trace
weights.
\fi
\ifnum\CRevisionRound>1
We derive the diffusion constant, exact odd-ring relaxation and the even-ring
period-two obstruction, including the one- and two-cell degeneracies.
Independent exact and symbolic checks audit the formulas. These source
results provide no target arithmetic or target-zero correspondence.
\fi
\end{abstract}
\noindent\textbf{Keywords:} multibaker map; primitive orbit; winding cocycle;
boundary coding; twisted determinant; deterministic diffusion.
\begin{otherlanguage}{chinese-simplified}
\begin{center}{\large 中文摘要}\end{center}
\begin{sloppypar}
本文在明确排除二进制边界的不变满测集上，建立有限环可逆多面包师映射
的完整原始周期轨道分类。每个几何周期点由混合二进制词唯一重建；长度
和位移决定提升后的最小周期、轨道条数、整数环绕数、反演配对及互为
倒数的稳定乘子。
\ifnum\CRevisionRound>0
有限维扭转输运行列式具有切比雪夫闭式。它的符号周期乘积含有两条不在
几何定义域中的齐次边界轨道；精确去除这两项得到全阶校正公式，并明确
区分不稳定乘子的逆权重与二维平坦迹权重。
\fi
\ifnum\CRevisionRound>1
进一步得到扩散常数、奇数环的精确松弛率、偶数环的二周期阻碍，以及
单格和双格退化情形。独立精确程序与符号计算核对公式。这些源系统定理
不构成目标算术结构或目标零点对应。
\fi
\end{sloppypar}
\noindent 关键词：多面包师映射；原始轨道；环绕余循环；边界编码；扭转行列式；确定性扩散。
\end{otherlanguage}

\section{The phase space and the endpoint problem}
The classical multibaker is a deterministic area-preserving realization of
nearest-neighbor diffusion; the model and its transport interpretation are
established in the literature \cite{gaspard,wojcik}. The present contribution
is an explicitly normalized proof and independent computational audit. We
claim no priority for the classical diffusion mechanism. The point requiring
particular care is that symbolic completeness and geometric completeness
depend on the chosen boundary convention.

Let $\mathcal D$ denote the dyadic rationals in $[0,1]$, let
$X=\{x\in(0,1):x\notin\mathcal D\}$, and set
\[
M_L=(\ZZ/L\ZZ)\times X\times X,\qquad L\geq1.
\]
For $s=\lfloor2x\rfloor$ define $d_s=2s-1$ and
\begin{equation}\label{eq:map}
B(j,x,y)=(j+d_s\bmod L,\ 2x-s,\ (y+s)/2).
\end{equation}
Both cell length and time step equal one. Normalized cell counting measure
times Lebesgue area gives an invariant probability measure on $M_L$.
Removing dyadic coordinates removes a null set; it does affect the exact
periodic-orbit convention, which is why it is stated before any census.

On a half-open unit square the all-one periodic code reconstructs the
excluded point $(1,1)$. Keeping that code while using a single-valued
half-open baker therefore overcounts a geometric orbit. Our symmetric
non-dyadic domain excludes both homogeneous codes and retains every mixed
periodic code. This makes the reversal exact on the stated domain.

\begin{lemma}[Coding and reversal]\label{lem:code}
The map $B$ is bijectively conjugate to the two-sided binary shift with cell
cocycle $d_s$, restricted to sequences whose two tails are not eventually
constant. It preserves area on each branch. The involution
$I(j,x,y)=(j,1-y,1-x)$ satisfies $IBI=B^{-1}$ and reverses the area form.
\end{lemma}
\begin{proof}
Use the unique expansions $x=.s_0s_1\cdots$ and
$y=.s_{-1}s_{-2}\cdots$. Formula \eqref{eq:map} shifts the symbolic sequence
and adds $d_{s_0}$ to the cell. Non-dyadic expansions are unique and neither
tail is eventually constant, so this coding and its inverse are defined
without endpoint identifications. With $t=\lfloor2y\rfloor$ one computes
\[
B^{-1}(j,x,y)=(j-d_t,\ (x+t)/2,\ 2y-t).
\]
Since $y\neq1/2$, the first symbol of $1-y$ is $1-t$; substitution proves
$IBI=B^{-1}$. The derivative on each branch is
$\operatorname{diag}(2,1/2)$, while $I^*(\dd x\wedge\dd y)=-\dd x\wedge\dd y$.
The images of the two branches partition the full-measure domain, establishing
measure preservation as well as the local Jacobian statement.
\end{proof}

\section{All geometric periods, powers and winding}
For a word $w=s_0\cdots s_{n-1}$, write $a(w)$ for its binary integer,
$w^{\rm rev}$ for its reversal, and $S(w)=\sum_t(2s_t-1)$. A mixed word
contains both symbols. A necklace is a word modulo cyclic rotation; it is
primitive when its least word period equals its length.

\begin{theorem}[Complete geometric atlas]\label{thm:atlas}
Every length-$n$ mixed word satisfying $L\mid S(w)$ gives, for each cell $j$,
exactly one point fixed by $B^n$, with
\begin{equation}\label{eq:coords}
x=\frac{a(w)}{2^n-1},\qquad y=\frac{a(w^{\rm rev})}{2^n-1}.
\end{equation}
These are all fixed points. A primitive mixed necklace of length $d$ and
displacement $S$ lifts to exactly $g=\gcd(L,S)$ geometric cycles, each with
\begin{equation}\label{eq:lift}
q=dL/g,\qquad W=S/g.
\end{equation}
Here $q$ is the least geometric period, $W$ is winding, and
$\gcd(L,0)=L$. Its $r$-fold repetition has time $rq$, winding $rW$,
multipliers $2^{rq},2^{-rq}$ and phase $\ee^{\ii r\phi W}$.
\end{theorem}
\begin{proof}
Affine iteration gives $x_n=2^nx-a(w)$ and
$y_n=2^{-n}y+2^{-n}a(w^{\rm rev})$. The fixed equations force
\eqref{eq:coords}. For mixed $w$, the resulting fractions lie strictly
between zero and one and have odd reduced denominators greater than one.
They are non-dyadic, and their periodic expansions realize the specified
word. The cell closes precisely when $L\mid S(w)$. Conversely every periodic
point has periodic coding by Lemma~\ref{lem:code}; the two homogeneous words
would give corners outside the domain.

For a primitive necklace, the internal coordinates first return after $d$
steps. At that time the cell translates by $S$. The order of this
translation is $h=L/g$, so the first full return occurs at $q=dh$. There
are $Ld$ choices of cell and word phase; partition into cycles of length
$dh$ gives $g$ distinct cycles. The total displacement is $hS$, yielding
$W=hS/L=S/g$. Each derivative is diagonal, giving the stated powers and
positive multipliers. No division by a reversal pairing has been made.
\end{proof}

The involution sends a cycle to the complement of its reversed word, up to
rotation. It sends $S,W$ to $-S,-W$ and keeps $q$ and multiplicity fixed.
Self-reversing cycles are allowed. The distinction between word period and
geometric period is substantial: the primitive word $001$ has $d=3$ and
$S=-1$, so on an $L$-cell ring it gives one cycle of period $3L$, not $L$
cycles of period three. The word $01$ has $S=0$, and gives $L$ cycles of
period two.

\begin{corollary}[All-period fixed census]
For every $L,n\geq1$,
\begin{equation}\label{eq:census}
F_{L,n}=L\left(\sum_{\substack{0\leq k\leq n\\L\mid(2k-n)}}
\binom nk-2\mathbf1_{L\mid n}\right).
\end{equation}
The number of geometric primitive cycles of period $q$ is
$q^{-1}\sum_{r\mid q}\mu(r)F_{L,q/r}$.
\end{corollary}
\begin{proof}
A length-$n$ word with $k$ ones has displacement $2k-n$ and there are
$\binom nk$ such words. Each closing word permits all $L$ cells. The two
homogeneous words close precisely when $L\mid n$ and must be subtracted.
Every period-$n$ fixed point belongs to one unique primitive cycle whose
period divides $n$; Möbius inversion gives the last assertion.
\end{proof}

\ifnum\CRevisionRound>0
\section{The tilted transport determinant and geometric correction}
For real $\phi$, define the finite backward transport operator
\begin{equation}\label{eq:transport}
(P_\phi f)_j=\tfrac12\ee^{\ii\phi/L}f_{j+1}
             +\tfrac12\ee^{-\ii\phi/L}f_{j-1}.
\end{equation}
The two labelled transitions are added when their destination cells coincide.
In particular, this definition does not discard either channel for $L=1,2$.
We consider $D_L(z,\phi)=\det(I-zP_\phi)$. This is a finite source transport
determinant with a specified observable space.

\begin{theorem}[Determinant and complete trace]\label{thm:det}
The full spectrum, with algebraic multiplicities, is
\[
\lambda_k(\phi)=\cos\frac{2\pi k+\phi}{L},\qquad0\leq k<L.
\]
Writing $T_L$ for the Chebyshev polynomial,
\begin{equation}\label{eq:cheb}
D_L(z,\phi)=2^{1-L}z^L\bigl(T_L(1/z)-\cos\phi\bigr)
=\prod_{k=0}^{L-1}(1-z\lambda_k(\phi)).
\end{equation}
The displayed Laurent expression is understood after cancellation as a
polynomial. For every $n\geq1$,
\begin{equation}\label{eq:trace}
\tr P_\phi^n=L2^{-n}\sum_{\substack{w\in\{0,1\}^n\\L\mid S(w)}}
\ee^{\ii\phi S(w)/L}.
\end{equation}
\end{theorem}
\begin{proof}
Fourier vectors $f_j=\ee^{2\pi\ii kj/L}$ diagonalize
\eqref{eq:transport}, giving the spectrum. For generic $\phi$, these are
distinct roots of $T_L(t)-\cos\phi$, since $T_L(\cos\theta)=\cos L\theta$.
The leading coefficient is $2^{L-1}$, including $L=1$. Equality of the
polynomials for generic $\phi$ extends by continuity to every real twist,
including multiple roots. Substitution $t=1/z$ proves \eqref{eq:cheb}.
Expanding the matrix power as labelled steps proves \eqref{eq:trace}:
diagonal entries require cell closure, and the accumulated phase is
$\ee^{\ii\phi S/L}$, not $\ee^{\ii\phi S}$.
\end{proof}

Eigenvalues equal to zero do not produce finite determinant roots. Each
nonzero eigenvalue contributes $z=1/\lambda_k$ with its multiplicity, so the
formula also specifies every degree drop and every source root count.
Replacing $\phi$ by $\phi+2\pi$ permutes the eigenvalue list.

\begin{theorem}[Boundary-corrected geometric product]\label{thm:zeta}
Let $p$ run over the geometric primitive cycles of Theorem~\ref{thm:atlas}.
For real $\phi$ and $|z|<1$, the absolutely convergent product is
\begin{equation}\label{eq:zeta}
\prod_p\left(1-(z/2)^{q_p}\ee^{\ii\phi W_p}\right)^{-1}
=\frac{(1-(z/2)^L\ee^{\ii\phi})(1-(z/2)^L\ee^{-\ii\phi})}
{D_L(z,\phi)}.
\end{equation}
The quotient gives its rational continuation, retaining cancellations.
\end{theorem}
\begin{proof}
Since $|\tr P_\phi^n|\leq L$, the series
$\sum_{n\geq1}z^n\tr P_\phi^n/n$ converges absolutely for $|z|<1$ and equals
$-\log D_L$. Group each closed labelled walk by its unique primitive cycle
and repetition. This yields the reciprocal Euler product over symbolic
lattice cycles, with inverse-unstable weight $2^{-q}$ and phase
$\ee^{\ii\phi W}$. Its only non-geometric cycles are the all-zero and all-one
cycles. Each has least lattice period $L$, and their windings are $-1,+1$.
Multiplying by the two corresponding factors removes them, proving
\eqref{eq:zeta}. They remain distinct labelled cycles even for $L=1,2$.
Alternatively the bound $F_{L,n}\leq L2^n$ controls the geometric logarithmic
series directly. Both arguments justify regrouping without any cancellation
being replaced by an absolute weight.
\end{proof}

There is an additional weight boundary. The inverse unstable multiplier is
$2^{-n}$, but the two-dimensional fixed-point denominator is
\begin{equation}\label{eq:flat}
|\det(I-DB^n)|=(2^n-1)^2/2^n.
\end{equation}
Its inverse is different. The finite transport determinant therefore has not
been proved to be a Fredholm determinant of the full phase-space Perron
operator. Neither an infinite-dimensional trace-class claim nor a target
divisor is licensed by \eqref{eq:cheb}.
\fi

\ifnum\CRevisionRound>1
\section{Hydrodynamic branch and exact ring relaxation}
\begin{theorem}[Diffusion and parity-sensitive rates]\label{thm:diffusion}
On the integer-cell lift with invariant internal Lebesgue data, the
displacement $X_n$ obeys
\[
\mathbb E X_n=0,\qquad \operatorname{Var}X_n=n,\qquad
\mathbb E\ee^{tX_n}=(\cosh t)^n,\qquad D=\tfrac12.
\]
The spatial Bloch branch is $\lambda(\kappa)=\cos\kappa$, and near zero
\[
\log\lambda(\kappa)=-\kappa^2/2-\kappa^4/12+O(\kappa^6).
\]
For odd $L\geq3$, the exact norm of $P_0^n$ on mean-zero cell observables is
$\cos(\pi/L)^n$. For even $L\geq2$, the eigenvalue $-1$ prevents convergence
to the uniform cell distribution at every integer time. For even $L\geq4$,
the two-step norm after removing both constant and parity modes is
$\cos^2(2\pi/L)$.
\end{theorem}
\begin{proof}
A prescribed future word of length $n$ is an interval of width $2^{-n}$ in
$x$, so future symbols are independent fair bits under the invariant measure.
The increments $2s_t-1$ are independent, centered and have variance one.
Multiplication of their moment-generating functions gives $(\cosh t)^n$,
and the diffusion convention is $D=\lim_n\operatorname{Var}(X_n)/(2n)$.
The Fourier characteristic function is $(\cos\kappa)^n$; Taylor expansion
gives the displayed local logarithm. The ring branch near zero twist is
$\cos(\phi/L)$, consistently using wave number $\kappa=\phi/L$.

At zero twist $P_0$ is real symmetric with eigenvalues
$\cos(2\pi k/L)$. For odd $L$, the largest absolute value outside $k=0$ is
$\cos(\pi/L)$; the spectral theorem makes the norm equality exact. For even
$L$, the mode $k=L/2$ is $(-1)^j$ and has eigenvalue $-1$. After removing
$k=0,L/2$, the largest squared eigenvalue is $\cos^2(2\pi/L)$. Squaring the
operator restricts motion to parity classes, giving the asserted two-step
relaxation factor on the mean-zero subspace of each parity class. No
probabilistic approximation is involved.
\end{proof}

\subsection{One cell, two cells, and the lazy control}
For $L=1$, $P_\phi=[\cos\phi]$ and $D_1=1-z\cos\phi$. The untilted cell
chain has only one invariant mode, so a nontrivial relaxation gap is not
defined. Its internal baker dynamics is still nontrivial: $F_{1,1}=0$,
$F_{1,2}=2$, and the primitive cycle $01$ survives.
For $L=2$, both labelled transitions reach the other cell and add to
$\cos(\phi/2)$; hence
\[
D_2=1-z^2\cos^2(\phi/2).
\]
The untilted chain alternates deterministically, with eigenvalues $1,-1$.
Each parity class is a singleton, so there is no remaining parity-class
relaxation mode. These facts do not remove either homogeneous symbolic cycle
from the correction in \eqref{eq:zeta}.

The usual nonabsolute gap for $L\geq3$ is $1-\cos(2\pi/L)$, which differs
from the odd-ring absolute gap. As a separate control, the lazy matrix
$Q=(I+P_0)/2$ has eigenvalues $\cos^2(\pi k/L)$ and gap
$\sin^2(\pi/L)$ for every $L\geq2$. This is an explicitly changed transport
process. It is not used to hide the original even-ring obstruction.

\section{Independent evidence and limits of the Route-A claim}
The all-period statements are proved above. The finite evidence contains
96 fixed-period rows, 188 winding-refined primitive rows, 2,296 necklace
lift rows and 326 direct rational geometric witnesses. Fixed and primitive
cutoffs are $n\leq12$, $L\leq8$; direct geometry uses $n\leq6$, $L\leq6$.
The producer uses binomial counts and primitive-necklace enumeration. The
checker independently uses signed-walk recursion, primitive-power inversion,
direct inverse-map and reversal tests, and Newton/exponential series.
It imports no producer code.

A third lane computes symbolic determinants of the actual tilted matrices,
checks 18 exact identities, and tests 1,200 Fourier cells at 80 decimal
digits with residual below $10^{-70}$. Two isolated working directories
reproduce the evidence byte for byte. Twenty-four repaired-hash mathematical
and metadata mutations and ten serialization attacks are rejected; three
smoke tests audit endpoints, parity and code independence. Strict JSON and
YAML reject ambiguous structure and type substitutions. No finite regression
is presented as proof of an infinite orbit census.

The phase-space geometry supports a formal quantization clue: the branches
are symplectic and $I$ is anti-symplectic. Quantum multibaker constructions
have a primary literature owner \cite{wojcik}; this paper constructs no
quantum operator or domain and claims no quantum orbit theorem. The integers
$L$, word lengths and windings carry no intrinsic rational-prime data.
The same identities hold for prime and composite ring sizes; parity governs
the transport obstruction. Consequently the strict tuple is
\[
(A0_{\rm FAIL},A1_{\rm WEAK},A2_{\rm FAIL},A3_{\rm FAIL},
 A4_{\rm FORMAL\_HINT}).
\]
The source atlas and source determinant are fully proved, while target A1
arithmetic carrying and target A2--A3 requirements remain unmet. The scope is
\texttt{NO\_BAD\_EULER\_OR\_ROOT\_NUMBER}; all target-claim flags and Route B
are false.
\fi

\section{Conclusion and revision record}
\ifcase\CRevisionRound
Round zero establishes the complete boundary-aware geometric orbit atlas,
including least periods after the cell lift, powers, winding and reversal.
Its exact census uses only mixed codes on the stated invariant domain.
\or
Round one adds the twisted transport determinant, its complete source
spectrum and the geometric correction for the two omitted homogeneous
cycles. It also separates the transport weight from the full flat-trace
denominator, closing the principal determinant-convention ambiguity.
\else
Round two adds the exact diffusion law, parity-sensitive relaxation theorem,
one- and two-cell boundary cases, and the independent evidence audit.
The resulting source theorem is complete under its declared domain and
weight conventions. An arithmetic carrier would be a new research
requirement, rather than a consequence of the transport identities.
\fi

\begin{thebibliography}{9}
\bibitem{gaspard} P. Gaspard, What is the role of chaotic scattering in
irreversible processes?, \emph{Chaos} \textbf{3} (1993), 427--442.
\href{https://doi.org/10.1063/1.165950}{doi:10.1063/1.165950}.
\bibitem{wojcik} D. K. W\'ojcik and J. R. Dorfman, Quantum multibaker maps:
Extreme quantum regime, \emph{Physical Review E} \textbf{66} (2002), 036110.
\href{https://arxiv.org/abs/cond-mat/0203494}{arXiv:cond-mat/0203494};
\href{https://doi.org/10.1103/PhysRevE.66.036110}{doi:10.1103/PhysRevE.66.036110}.
\end{thebibliography}
\end{document}
